\documentclass[11pt,a4paper]{article}
\usepackage[utf8]{inputenc}
\usepackage[T1]{fontenc}
\usepackage[spanish,es-nodecimaldot]{babel}
\usepackage{lmodern,amsmath,amssymb,booktabs,tabularx,array}
\usepackage[margin=2.1cm]{geometry}
\usepackage{microtype,xcolor}
\usepackage[colorlinks=true,linkcolor=blue!45!black,urlcolor=blue!45!black,hypertexnames=false]{hyperref}
\newcommand{\Rminus}{R^{-}}
\newcommand{\reserve}{\underline b}
\newcommand{\code}[1]{\texttt{\small #1}}
\setlength{\parindent}{0pt}
\setlength{\parskip}{6pt}
\setlength{\tabcolsep}{5pt}
\renewcommand{\arraystretch}{1.2}
\title{\vspace{-1cm}Modelo de optimización de ruteo eléctrico\\
\large Formulación base y extensiones operacionales configurables}
\author{Proyecto AWS Breaking Barriers Chile 2026 --- CASE / FCFM}
\date{Actualización: 15 de octubre de 2026}
\begin{document}
\maketitle

\textbf{Objetivo.} Maximizar entregas factibles y, entre soluciones con igual
cumplimiento, minimizar consumo neto efectivo de batería. El tiempo de cómputo
se reporta; no es una expresión optimizada dentro del MILP.

\textbf{Formulación vigente sin $H$.} El backend predeterminado
\code{scip-routes} implementa (25)--(29), en la sección~7.
Las ecuaciones posicionales (1)--(24) se conservan como referencia Pyomo.
Sus condiciones físicas y operacionales se mantienen en las rutas factibles,
pero no su cota artificial de movimientos.

\textbf{Representaciones.} La base (1)--(16) describe calles elementales
dirigidas, admite regeneración y cuenta movimientos viales. En el flujo sin
$H$, \code{service-graph} usa Dijkstra si los consumos son no negativos y no hay
plazos; en los demás casos compone perfiles de batería y conserva alternativas
temporales e intersecciones cuando es necesario (sección~8).
La referencia posicional conserva caminos fijos con $H$ en traslados entre
servicios: no equivale al modelo con $H$ vial. Las extensiones (17)--(24)
se añaden a la misma plantilla.

\section{Parámetros y conjuntos}
\begin{tabularx}{\textwidth}{@{}l X@{}}
\toprule
Símbolo & Significado\\\midrule
$G=(V,A)$ & Grafo dirigido; $0\in V$ es el depósito. Se añade $(0,0)$ con energía y tiempo cero.\\
$C,F\subseteq V$ & Clientes y cargadores. En los CSV son tipos de nodo distintos; $0\notin C\cup F$.\\
$K$ & Vehículos disponibles, homogéneos. Escenario oficial: dos Nissan Leaf 2013.\\
$F_k\subseteq F$ & Estaciones físicamente compatibles y habilitadas para $k$.\\
$B>0$ & Capacidad útil, kWh; es configurable, no una medición del dataset.\\
$\reserve$ & Reserva durante todo el recorrido, $0<\reserve\le b^{ini}$.\\
$b^{ini}$ & Energía inicial: $0.6B$ en el escenario oficial.\\
$e_{ij}$ & Energía nominal del arco, kWh: consumo positivo y regeneración potencial negativa.\\
$H\in\mathbb N_{\ge1}$ & Máximo de movimientos por vehículo; limita la representación, no la jornada.\\
$R,\Rminus$ & Posiciones $R=\{0,\ldots,H\}$ y posiciones con salida $\Rminus=\{0,\ldots,H-1\}$.\\
$M_B$ & Cota de saturación $B+\max_{(i,j)\in A}|e_{ij}|$; con margen, usar $\widetilde e$ en su lugar.\\
\bottomrule
\end{tabularx}

Los arcos elementales no pueden ocultar una descarga seguida de regeneración.
En la reducción no negativa de Dijkstra, comprobar la reserva al final del
camino basta para sus segmentos. Con regeneración se necesitan los perfiles
(30)--(31), incluida la reserva interior. La auditoría siempre expande y
comprueba las calles originales.

\section{Variables y dominios}
\begin{tabularx}{\textwidth}{@{}l X@{}}
\toprule
Variable & Significado y dominio\\\midrule
$y_i\in\{0,1\}$ & Entrega cumplida a $i\in C$.\\
$z_{ikr}\in\{0,1\}$ & Ubicación de $k$ en $i\in V$, posición $r\in R$.\\
$x_{ijkr}\in\{0,1\}$ & Uso del arco $(i,j)\in A$, $r\in\Rminus$.\\
$a_{ikr}\in\{0,1\}$ & Entrega a $i\in C$ en la posición $r\in\Rminus$. No se activa por pasar por el nodo.\\
$q_{kr}\ge0$ & Energía que entra a la batería durante la parada, $r\in\Rminus$.\\
$b_{kr}\in[\reserve,B]$ & Energía al llegar, $r\in R$.\\
$\bar b_{kr}\in[\reserve,B]$ & Energía al salir, $r\in\Rminus$.\\
$u_{kr},E_{kr}\in\mathbb R$ & Batería previa a saturación y energía efectiva del movimiento.\\
$\gamma_{kr}\in\{0,1\}$ & Indicador de saturación, $r\in\Rminus$.\\
$E_{\mathrm{total}}\in\mathbb R$ & Consumo neto efectivo total.\\
\bottomrule
\end{tabularx}

\section{Formulación base}
Las restricciones se imponen para todos los índices indicados.
\begin{align}
\sum_{i\in V} z_{ikr}&=1 && k\in K,\ r\in R, \tag{1}\\
z_{0k0}=z_{0kH}&=1 && k\in K, \tag{2}\\
\sum_{j:(i,j)\in A}x_{ijkr}&=z_{ikr} && i\in V,\ k\in K,\ r\in\Rminus, \tag{3}\\
\sum_{i:(i,j)\in A}x_{ijkr}&=z_{jk,r+1} && j\in V,\ k\in K,\ r\in\Rminus, \tag{4}\\
a_{ikr}&\le z_{ikr} && i\in C,\ k\in K,\ r\in\Rminus, \tag{5}\\
\sum_{k\in K}\sum_{r\in\Rminus}a_{ikr}&=y_i && i\in C, \tag{6}\\
b_{k0}=b^{ini},\quad q_{k0}&=0 && k\in K, \tag{7}\\
\bar b_{kr}&=b_{kr}+q_{kr} && k\in K,\ r\in\Rminus, \tag{8}\\
0\le q_{kr}&\le B-b_{kr} && k\in K,\ r\in\Rminus, \tag{9}\\
q_{kr}&\le B\sum_{f\in F_k}z_{fkr} && k\in K,\ r\in\Rminus. \tag{10}
\end{align}
La continuidad permite revisitar calles, cargadores y depósito. No se impone
permanencia tras un retorno intermedio. El arco $(0,0)$ permite inactividad
y relleno de posiciones.
\begin{align}
u_{kr}&=\bar b_{kr}-\sum_{(i,j)\in A}e_{ij}x_{ijkr}, \tag{11}\\
b_{k,r+1}&\le u_{kr},\qquad b_{k,r+1}\le B, \tag{12}\\
b_{k,r+1}&\ge u_{kr}-M_B\gamma_{kr},\qquad
b_{k,r+1}\ge B-M_B(1-\gamma_{kr}), \tag{13}\\
E_{kr}&=\bar b_{kr}-b_{k,r+1}, \tag{14}\\
E_{\mathrm{total}}&=\sum_{k\in K}\sum_{r\in\Rminus}E_{kr}. \tag{15}
\end{align}
Las ecuaciones (11)--(14) rigen para $k\in K,r\in\Rminus$.
(12)--(13) representan exactamente $b_{k,r+1}=\min(B,u_{kr})$.
La regeneración que excede capacidad no se almacena ni se acredita.
Los dominios conservan reserva también al regreso.

El problema base es:
\begin{equation}
\boxed{\operatorname{lexmin}\left(-\sum_{i\in C}y_i,\ E_{\mathrm{total}}\right)
\quad\text{sujeto a dominios y (1)--(15).}} \tag{16}
\end{equation}
Se resuelven dos MILP: maximizar entregas y luego fijar el mejor cumplimiento
encontrado para minimizar energía. Si la primera fase no prueba su máximo,
la segunda es una mejora condicionada, no un óptimo lexicográfico demostrado.
Se conservan cotas, gaps, estados y tiempos \emph{por fase}.

\section{Extensiones opcionales sobre los caminos fijos}
Los CSV no entregan tiempos de servicio, jornadas ni ventanas de recepción.
Tampoco acreditan potencia efectiva aceptada por el vehículo, eficiencia o
disponibilidad en tiempo real de los cargadores. Estos valores son entradas
del cliente o supuestos explícitos del escenario, no reglas oficiales.

\subsection{Margen energético, estaciones y reserva final}
Sean $\rho\ge0$ el margen de consumo y $D\subseteq F$ las estaciones
deshabilitadas. Se utiliza:
\begin{equation}
\widetilde e_{ij}=e_{ij}+\rho\max(e_{ij},0),\qquad
F_k^{\mathrm{activo}}=F_k\setminus D. \tag{17}
\end{equation}
El máximo es sobre un parámetro, no sobre una variable. Se sustituye $e$
por $\widetilde e$ en (11) y en $M_B$, y $F_k$ por
$F_k^{\mathrm{activo}}$ en (10). La energía nominal permanece en el snapshot.
El margen se aplica una sola vez, también en el validador y la expansión vial.
En la variante no negativa equivale a multiplicar toda energía por $1+\rho$,
sin cambiar los caminos de Dijkstra. No es una garantía probabilística.

Si se configura una fracción final $\beta\in[0,1]$:
\begin{equation}
b_{kH}\ge\beta B\qquad k\in K. \tag{18}
\end{equation}
La reserva original $\reserve$ sigue vigente. (18) no se impone en posiciones
intermedias. Si $\beta B>b^{ini}$, incluso un vehículo sin entregas necesita
un recorrido factible con recarga; no se debilita (18) para permitir inactividad.
Deshabilitar una estación prohíbe recargar allí, no atravesar su ubicación.

\subsection{Viaje, servicio y recarga}
Todos los tiempos se expresan en segundos desde una salida común en $t=0$.
$\tau_{ij}\ge0$ es el tiempo del camino seleccionado y $\tau_{00}=0$.
Se suman los tiempos explícitos de sus calles; donde faltan, solo se permite
una velocidad asumida $v>0$ declarada:
$\tau_{\mathrm{calle}}=3.6\,d_{\mathrm{calle}}/v$, con distancia en metros y
$v$ en km/h. No se interpreta el límite legal como velocidad real.
$s_i\ge0$ es el tiempo de servicio, común o particular por cliente.

Para incluir tiempo de recarga se añaden variables continuas $q_{fkr}\ge0$,
solo para $f\in F_k^{\mathrm{activo}}, k\in K,r\in\Rminus$:
\begin{equation}
q_{kr}=\sum_{f\in F_k^{\mathrm{activo}}}q_{fkr},\qquad
0\le q_{fkr}\le Bz_{fkr}. \tag{19}
\end{equation}
Sea $P_f>0$ la potencia efectiva compatible, del lado de la red, y
$0<\eta\le1$ su eficiencia constante. Esta versión usa la misma potencia
configurada por estación para los vehículos homogéneos. No se toma
\code{max\_power\_kw} como potencia por conector ni como potencia aceptada.
Se calculan expresiones lineales:
\begin{equation}
S_{kr}^{\mathrm{dur}}=\sum_{i\in C}s_i a_{ikr},\quad
C_{kr}^{\mathrm{dur}}=\sum_{f\in F_k^{\mathrm{activo}}}
\frac{3600q_{fkr}}{\eta P_f},\quad
V_{kr}^{\mathrm{dur}}=\sum_{(i,j)\in A}\tau_{ij}x_{ijkr}. \tag{20}
\end{equation}
Cuando no se activa tiempo de recarga, $C_{kr}^{\mathrm{dur}}=0$ y (19) no
se construye. Para cualquier escenario temporal con cargadores habilitados,
la configuración exige activarlo y proporcionar las potencias efectivas.
Así no se permiten recargas instantáneas dentro de una jornada.

Sin ventanas no se necesitan relojes ni espera. Si se configura una duración
máxima $T>0$:
\begin{equation}
\sum_{r\in\Rminus}
\left(V_{kr}^{\mathrm{dur}}+S_{kr}^{\mathrm{dur}}+C_{kr}^{\mathrm{dur}}\right)
\le T \qquad k\in K. \tag{21}
\end{equation}
El límite es común a la flota actual. El tiempo de servicio aislado, sin
plazo ni ventanas, afecta al reporte pero no a los objetivos de (16).

\subsection{Ventanas de recepción completas y espera}
Al activarse ventanas se exige también un $T$ finito. Sea $C_W\subseteq C$
el subconjunto configurado y $[L_i,U_i]$ su ventana, con $0\le L_i\le U_i$.
Se añaden llegadas $t_{kr}\in[0,T]$, $r\in R$, e inicios de actividad
$S_{kr}\in[0,T]$, $r\in\Rminus$:
\begin{equation}
\begin{aligned}
t_{k0}&=0,\\
S_{kr}&\ge t_{kr},\\
t_{k,r+1}&=S_{kr}+S_{kr}^{\mathrm{dur}}
+C_{kr}^{\mathrm{dur}}+V_{kr}^{\mathrm{dur}}
&& k\in K,\ r\in\Rminus.
\end{aligned} \tag{22}
\end{equation}
La espera es $S_{kr}-t_{kr}$. No hay un consumo auxiliar de espera o servicio
en esta extensión; permanece como supuesto explícito.
\begin{equation}
\begin{aligned}
S_{kr}&\ge L_i a_{ikr},\\
S_{kr}+s_i a_{ikr}
&\le U_i+\max(0,T-U_i)(1-a_{ikr})
&& i\in C_W,\ k\in K,\ r\in\Rminus.
\end{aligned} \tag{23}
\end{equation}
Si $a_{ikr}=1$, la entrega empieza después de abrir y termina antes de
cerrar. Si $a_{ikr}=0$, las cotas $0\le S_{kr}\le T$ hacen ambas condiciones
redundantes. Estos coeficientes temporales no reutilizan $M_B$. Los clientes
sin ventana permanecen sujetos a la jornada, pero no a (23).
Pasar por un cliente, incluso dentro de un camino agregado, no activa servicio.

Con las opciones habilitadas, el problema resuelto es:
\begin{equation}
\boxed{\operatorname{lexmin}\left(-\sum_{i\in C}y_i,\ E_{\mathrm{total}}\right)
\quad\text{sujeto a (1)--(15) y las extensiones activas (17)--(23).}} \tag{24}
\end{equation}
No se añade un objetivo de puntualidad, duración ni costo que altere la
prioridad oficial. Se reporta el calendario más temprano de las mismas
rutas y recargas; bajo los supuestos constantes de esta sección, quitar
espera sobrante no altera energía ni entregas y no empeora factibilidad.

\section{Configuración e interpretación}
\small
\begin{tabularx}{\textwidth}{@{}>{\ttfamily}l X@{}}
\toprule
\normalfont Campo de \code{operations} & Desactivación y significado\\\midrule
service\_seconds & 0: sin servicio. Segundos comunes por entrega.\\
service\_seconds\_by\_customer & \{\}: sin excepciones. Un valor por cliente reemplaza al común.\\
max\_route\_seconds & null: sin jornada máxima adicional.\\
assumed\_speed\_kph & null: exigir tiempos explícitos si se activa el bloque temporal.\\
include\_charging\_time & false: sin bloque temporal de recarga; condicionado como se indica antes.\\
charging\_power\_kw & null: sin potencia común; aportar potencias por estación si se necesita cargar.\\
charging\_power\_kw\_by\_station & \{\}: usar potencia común. Un valor particular la reemplaza.\\
charging\_efficiency & 1: hipótesis de recarga sin pérdidas. Se configura en $(0,1]$.\\
final\_reserve\_soc & null: solo la reserva de la base.\\
energy\_margin\_fraction & 0: energía nominal.\\
disabled\_charger\_ids & []: no excluir estaciones adicionales.\\
use\_customer\_time\_windows & false: no construir relojes ni ventanas.\\
customer\_time\_windows & \{\}: sin ventanas. Cada entrada contiene open\_seconds y close\_seconds.\\
\bottomrule
\end{tabularx}
\normalsize

Todos los valores por defecto son neutros. Se validan referencias, unidades,
finitud, dependencias y potencias configuradas. Un booleano no sustituye datos
faltantes. El ejemplo de configuración usa valores ilustrativos, no parámetros
medidos ni exigencias del PPTX.

\section{H y arranque de la referencia posicional}
La prueba existente de $H=|C|+1$ necesita energías no negativas, desigualdad
triangular dirigida y un testigo completo cuya energía cabe entre batería
inicial y reserva, además de ausencia de restricciones que impidan eliminar
paradas. Una jornada, ventanas o reserva final mayor que la reserva base
desactivan esta prueba y su cota estructural en la implementación actual.
\textbf{Esto no prueba que $|C|+1$ sea insuficiente}: significa que el programa
no tiene un certificado válido para esa extensión.

En esos casos, \code{auto} se detiene y solicita un $H$ explícito. No se
construye un MILP enorme automáticamente. Se conserva el alcance condicionado
al $H$ elegido. El margen energético y las estaciones deshabilitadas se
incorporan al análisis cuando sus restantes hipótesis siguen siendo válidas.

Una heurística greedy acotada puede repartir entregas entre vehículos sin
recargar y producir un incumbente parcial; no certifica $H$ ni optimalidad.
Se comprueban batería, retorno, jornada, ventanas y las calles expandidas antes
de enviar cualquier arranque. No se fija ninguna variable del solver. Las
huellas del snapshot incluyen todas las opciones: cambiar de escenario invalida
el arranque y los hashes anteriores.

La auditoría independiente reproduce las transiciones y los intervalos
reportados. La duración operacional es el máximo de las duraciones de los
vehículos, con salida común; su suma se reporta por separado como tiempo de
vehículos. La electricidad comprada es $\sum q/\eta$ solo cuando la eficiencia
ha sido configurada en el bloque correspondiente.

No se modelan colas, puertos ocupados, apertura de estaciones, carga no lineal,
capacidad de paquetes, reabastecimiento, tarifa eléctrica, climatización
durante paradas ni tráfico variable con la hora. El depósito sigue sin recarga.
Los caminos de energía mínima pueden ser más lentos que otras alternativas:
las ventanas y jornadas se optimizan sobre los caminos fijos disponibles,
sin afirmar equivalencia con la red vial completa.

\clearpage
\section{Reformulación vigente por rutas sin $H$}
Sea $\mathcal R_k$ el conjunto de recorridos finitos físicamente factibles del
vehículo $k$, sin una cota de longitud. Cada recorrido parte y termina en el
depósito, puede revisitar calles, cargadores y depósito, y distingue tránsito
de servicio. Conserva reserva, capacidad, recarga parcial continua, compatibilidad,
saturación de regeneración y las opciones (17)--(24).
La ruta inactiva solo se incluye si satisface también la reserva final.

Para $r\in\mathcal R_k$, sean $a_{ir}\in\{0,1\}$ el indicador de entrega,
$Q_r$ la carga acumulada y $c_r=b^{ini}+Q_r-b_r^{final}$ la energía neta efectiva.
La asignación $w_{ki}$ es binaria; los pesos de rutas $\lambda_{kr}$ son continuos.
\begin{align}
\sum_{r\in\mathcal R_k}\lambda_{kr}&=1 && k\in K, \tag{25}\\
\sum_{r\in\mathcal R_k}a_{ir}\lambda_{kr}&=w_{ki} && k\in K,\ i\in C, \tag{26}\\
\sum_{k\in K}w_{ki}&\le1,\quad w_{ki}\in\{0,1\},\quad\lambda_{kr}\ge0, \tag{27}\\
\operatorname{lexmin}\Bigl(-\sum_{ki}w_{ki},&
\ \sum_{kr}c_r\lambda_{kr}\Bigr). \tag{28}
\end{align}
Se ramifica sobre $w$, no sobre columnas desconocidas. Si $w_k$ es entero,
una combinación convexa de vectores binarios $a_r$ solo puede mezclar vectores
idénticos. Por tanto todas las columnas positivas de $k$ atienden el mismo
subconjunto; seleccionar la más barata recupera una ruta entera físicamente
factible sin empeorar el objetivo. Esto demuestra equivalencia entera.

El maestro completo convexifica la proyección local de cada vehículo sobre
entregas y coste. \textbf{No demuestra idealidad global}: el acoplamiento entre
vehículos puede originar puntos fraccionarios. Para una terna $S\subseteq C$ se
puede reforzar la relajación con el corte válido:
\begin{equation}
\sum_{kr}\left\lfloor\frac{\sum_{i\in S}a_{ir}}{2}\right\rfloor\lambda_{kr}
\le1,\qquad |S|=3. \tag{29}
\end{equation}
Toda columna futura recibe su coeficiente en (29), y el pricing incorpora
el multiplicador dual correspondiente.

\subsection{Pricing continuo y criterio de parada}
Las etiquetas conservan nodo, conjunto atendido y un poliedro de recursos
$(b,t,Q)$. El reloj se necesita en el pricing cuando existe una jornada física.
Las cargas cubren $0\le q\le B-b$ sin discretización de SoC. Las regiones de
saturación y espera se separan antes de usar transformaciones afines y cascos
convexos. Se guardan testigos de cargas para reproducir cada ruta.

Solo se descarta un estado cuando otro, en el mismo nodo y con el mismo
conjunto atendido, tiene al menos tanta batería, no más tiempo y no más carga
acumulada. Esta última componente es necesaria para comparar correctamente
energía efectiva con regeneración. No se impone máximo de etiquetas ni visitas.
La enumeración puede ser costosa; un timeout no prueba pricing completo.

SCIP ejecuta pricing de coste reducido y Farkas. Si el oráculo se interrumpe,
el LP restringido no acredita una cota global. Se conservan cotas analíticas:
entregas $\le |C|$, energía $\ge0$ con costes no negativos y
energía $\ge |K|(b^{ini}-B)$ en general. Una solución se valida independientemente;
la optimalidad exige cotas válidas que cierren ambas fases.

\subsection{Reducción certificada a un tour}
Con costes no negativos, métrica dirigida, ausencia de jornada, ventanas y
reserva final adicional, supóngase un tour completo de coste
$U\le b^{ini}-\reserve$. Un óptimo de flota cuesta a lo sumo $U$.
Concatenar sus rutas y quitar paradas sin entrega no aumenta el consumo por
desigualdad triangular. El tour obtenido cabe en la batería inicial.
Existe por tanto un óptimo sin recargas con un solo vehículo activo.

La implementación verifica esas hipótesis antes de usar un modelo dirigido
de grados y cortes DFJ de subtours. No usa posiciones, MTZ ni Big-M de batería.
El testigo heurístico no queda probado óptimo; la resolución DFJ determina
coste y cota. Esta reducción conserva el óptimo, no todas las rutas factibles.

\subsection{Incumbentes y relajación energética}
Se generan incumbentes mediante inserción, búsqueda local, partición de un
orden de entregas y una vecindad de recarga parcial. Cada resultado se valida
independientemente. Estas heurísticas no restringen $\mathcal R_k$: el pricing
exacto conserva múltiples recargas, revisitas y recursos continuos.
Un maestro con un conjunto finito de columnas conocidas sirve como heurística
primal; su cota nativa no certifica el maestro completo.

Con servicio completo y energías no negativas, sea $d_{ij}$ el coste de un
camino mínimo dirigido en el grafo de entrada. Concatenar los recorridos de
la flota en el depósito y atajar entre eventos de entrega usando $d$ produce
un tour cuyo coste no supera la energía de la flota. Así, una asignación
dirigida y las relajaciones con cortes de subtours proporcionan cotas inferiores
válidas al eliminar batería, recargas, tiempos y reparto entre vehículos.
La implementación resta un margen conservador de acumulación numérica.
Este argumento no requiere que los arcos directos ya satisfagan exactamente
la desigualdad triangular, pues se construye su clausura por caminos mínimos.
No se aplica con regeneración ni cuando quedan clientes sin atender.
Un tour relajado solo se incorpora como solución después de repararlo y validar
su factibilidad física. Se conserva la formulación (25)--(29).

\subsection{Orquestación del LLM y alcance}
Por autorización del usuario, Bedrock puede tomar decisiones en eventos
explícitos: antes de buscar, después del arranque, ante estancamiento y desde
callbacks de árbol, LP, pricing, separación e incumbentes.
La inferencia es asíncrona: otro hilo recibe solo JSON y los objetos nativos
permanecen en el hilo de SCIP. El límite predeterminado es cuatro llamadas y una
asignación de inferencia del 25\,\% del presupuesto global, ambos configurables.
La frecuencia está acotada; no se consulta remotamente por cada nodo.
Una decisión vencida no se aplica ni bloquea nuevas consultas presupuestadas.
No se presupone cancelada la solicitud remota: su uso queda desconocido hasta
recibir respuesta. Tiempo, tokens y cambios efectivos se registran por decisión.

El LLM controla heurísticas, operadores, reparto interno de tiempo, relajaciones,
pricing, cortes compatibles y parámetros nativos verificados. Los campos de
etapas terminadas quedan inactivos. No cambia datos, física, tolerancias,
presupuesto global ni criterios de prueba. Los callbacks ejecutan políticas
locales mientras avanza la inferencia; una propuesta inválida conserva la última política
válida y se registra como rechazo. El motor mantiene el pricing exacto ante
fallos de AWS.
El contexto incluye características y recursos de la instancia, duales, progreso,
rutas incumbentes y APIs implementadas. Se exportan el grafo físico y el MPS del
maestro restringido. Este último no contiene las columnas futuras ni representa
por sí solo el problema completo; cualquier recorte enviado se declara.
En snapshots antiguos, \code{movement\_limit} es metadato heredado y no interviene
en esta formulación. Las rutas resultantes tienen longitud variable.

El certificado sigue referido al grafo de entrada y al modelo energético.
La nueva preparación comprueba cuándo conserva el óptimo vial sin $H$; eso
no sustituye la prueba de optimalidad de ambas fases. Un catálogo heredado de
Dijkstra con plazos solo certifica sus caminos disponibles. Detalles y fuentes:
\code{docs/dantzig-wolfe.md} y \code{docs/resource-paths.md}.

\subsection{Búsqueda especializada y pricing por coste reducido}
ALNS destruye grupos aleatorios, relacionados, de alto desvío o de una ruta,
y repara mediante inserción greedy o regret. Una población memética combina
órdenes de entrega con elitismo y diversidad de aristas. Cada orden se repara
sobre el multigrafo físico, conservando alternativas de calles, recarga continua,
ventanas y revisitas. Un sub-SCIP combina columnas conocidas como heurística
primal; su cota no se transfiere al maestro completo.

La consulta exacta se reinicia con cada dual: sus podas nunca se reutilizan
con otro dual. Para una etiqueta en $v$ con clientes servidos $M$ y poliedro
$\mathcal P$ de batería $b$, tiempo y carga acumulada $Q$, una cota de la
energía final es
\[
\underline E=
\begin{cases}
b^{ini}+\min_{\mathcal P}(Q-b)+d(v,0),&\text{si no hay regeneración},\\
b^{ini}+\min_{\mathcal P}Q-B,&\text{en general}.
\end{cases}
\]
El primer caso usa una distancia mínima no negativa al depósito; el segundo
usa que la recarga acumulada no decrece y la batería final no supera $B$.
Para los premios duales $\pi_i$, una cota superior del premio de una continuación
es $\sum_{i\in M}\pi_i+\sum_{i\notin M}\max(0,\pi_i)$.
Para cada corte subset-row $S$, se usa su máximo coeficiente posible si su
dual es positivo y el coeficiente ya acumulado si es negativo.
Restar estos premios y el dual de convexidad de la cota energética produce
una cota inferior conservadora de coste reducido. En Farkas el peso de energía
es cero. Solo se poda si no puede existir una columna mejorante, incluyendo
un margen numérico conservador. No hay límites de movimientos ni etiquetas.

El branching actúa sobre asignaciones binarias de clientes, la selección de
nodos sigue políticas locales y el separador agrega triples subset-row como
restricciones modificables. Todas las columnas posteriores reciben sus
coeficientes. Las columnas complementarias de coste reducido no negativo
pueden admitirse como ayuda primal, pero no acreditan completitud del pricing.
La interrupción descarta la cota del LP restringido.

\clearpage
\section{Preprocesamiento de caminos con recursos}

El maestro (25)--(29) se mantiene. Se prepara por separado un multigrafo que
conserva depósito, clientes, cargadores y las intersecciones que no resulte
conveniente eliminar. Las conexiones tienen identificadores propios y testigos
de calles; el pricing selecciona entre ellas y propaga recursos continuos.
No se agrega un límite de movimientos, etiquetas, recargas o caminos por par.

\subsection{Perfil físico y composición}

Sea $P$ un camino sin actividades interiores y sean $e_1,\ldots,e_n$ sus consumos
después de aplicar (17) a cada descarga. Defínanse $S_0=0$,
$S_j=\sum_{\ell=1}^j e_\ell$, $m=\min_j S_j$, $M=\max_j S_j$ y
$D=\max_{0\le i\le j\le n}(S_j-S_i)$. El camino tiene algún estado inicial
factible si $M\le B-\reserve$ y $D\le B-\reserve$. Su perfil es
\begin{equation}
 L_P=\reserve+M,\quad E_P=S_n,\quad U_P=B-E_P+m,\qquad
 F_P(b)=\min\{b-E_P,U_P\},\quad b\in[L_P,B]. \tag{30}
\end{equation}
El consumo efectivo es $b-F_P(b)$; no se sustituye por la suma firmada $E_P$.
El dominio captura déficits interiores y $U_P$ captura saturación.
Para concatenar $P$ y $Q$, sin actividad en su unión:
\begin{equation}
\begin{split}
E_{PQ}&=E_P+E_Q,\\
L_{PQ}&=\max\{L_P,L_Q+E_P\},\\
U_{PQ}&=\min\{U_P-E_Q,U_Q\}.
\end{split} \tag{31}
\end{equation}
La concatenación es inviable si $U_P<L_Q$ o $L_{PQ}>B$.
Los tiempos y distancias se suman. El margen ya incorporado al perfil no se
aplica nuevamente en el pricing.

\subsection{Dominancia y eliminación parcial}

Entre conexiones con iguales extremos, $P$ domina $Q$ si
$L_P\le L_Q$, $F_P(b)\ge F_Q(b)$ para todo $b\in[L_Q,B]$ y, cuando hay jornada,
$T_P\le T_Q$. La diferencia entre dos funciones de (30) es monótona; basta
comprobar ambos extremos del intervalo. Las comparaciones de dominancia no
relajan tolerancias.

Eliminar una intersección sin actividades consiste en componer cada entrada
con cada salida. Todo recorrido por ella tiene esa representación y cada
conexión nueva se expande al mismo recorrido vial. Una conexión dominante
permite llegar antes y con más batería; puede esperarse y reducir una carga
parcial posterior. Por $c_r=b^{ini}+Q_r-b_r^{final}$, la sustitución no empeora
el consumo efectivo. Esto conserva el óptimo, aunque quite recorridos dominados.

Un autoarco dominado por permanecer en el nodo se elimina. Los ciclos que
recuperan energía se conservan y una intersección con un autoarco útil no se
elimina. Por tanto, no se supone ausencia de ciclos energéticos negativos.
Si una eliminación genera demasiadas combinaciones o se agota el presupuesto
de preparación, se conserva la intersección para que la explore el pricing.
Detener la compresión no descarta caminos pendientes.

\subsection{Alcance y auditoría}

La reducción preserva el óptimo vial sin $H$ bajo tiempos constantes, espera
sin consumo, recargas parciales y actividades solo en servicios. No cubre
tráfico dependiente de hora, consumo por carga transportada, objetivos de
distancia ni consumo auxiliar durante espera.
Con energías no negativas y sin plazos, Dijkstra sigue siendo suficiente por
monotonicidad de batería y objetivo energético.

Se comprueban perfiles contra prefijos viales y se expanden todas las calles
de la solución para validarlas independientemente. Ventanas, jornadas, recargas,
reserva final y clientes atendidos se comprueban sobre esa expansión.
La caché depende de la red, escenario, parámetros y versión. Las alternativas
y el pricing pueden crecer mucho: no hay una garantía de tiempo polinomial.
El presupuesto de preparación y el de resolución se registran por separado.
\end{document}
